% ---------------------------------------------------------------------------
% Preâmbulo: pacotes, aspeto e macros do relatório.
% ---------------------------------------------------------------------------

\usepackage[T1]{fontenc}
\usepackage[utf8]{inputenc}
\usepackage[portuguese]{babel}
\usepackage{lmodern}
\usepackage{microtype}
\usepackage[a4paper, margin=2.5cm, headheight=14pt]{geometry}
\usepackage{setspace}
\onehalfspacing

% Matemática
\usepackage{amsmath, amssymb, mathtools, bm}

% Tabelas
\usepackage{array, booktabs, tabularx, longtable, multirow, makecell}
\usepackage{colortbl}
\usepackage{threeparttable}
\usepackage{pdflscape}

% Figuras
\usepackage{graphicx}
\usepackage{float}
\usepackage[font=small, labelfont=bf, labelsep=period]{caption}
\usepackage{subcaption}
\usepackage{tikz}
\usetikzlibrary{arrows.meta, positioning, shapes.geometric, shapes.multipart,
                fit, calc, backgrounds, decorations.pathreplacing}
\usepackage{pgfplots}
\pgfplotsset{compat=1.16}

% Diversos
\usepackage{xcolor}
\usepackage{enumitem}
\usepackage{pifont}
\usepackage{csquotes}
\usepackage{url}
\usepackage{fancyhdr}
\usepackage{titlesec}
\usepackage[hidelinks, unicode, pdfusetitle]{hyperref}

% ---------------------------------------------------------------------------
% Cores (inspiradas na identidade do Iscte; ajustáveis)
% ---------------------------------------------------------------------------
\definecolor{iscte}{RGB}{0, 51, 102}
\definecolor{acento}{RGB}{0, 112, 192}
\definecolor{cinza}{RGB}{110, 110, 110}
\definecolor{fundotab}{RGB}{235, 241, 248}
\definecolor{bom}{RGB}{0, 128, 0}
\definecolor{mau}{RGB}{190, 30, 45}
\definecolor{parcial}{RGB}{200, 130, 0}

% ---------------------------------------------------------------------------
% Títulos de secção e cabeçalhos
% ---------------------------------------------------------------------------
\titleformat{\section}{\Large\bfseries\color{iscte}}{\thesection}{0.8em}{}
\titleformat{\subsection}{\large\bfseries\color{iscte}}{\thesubsection}{0.8em}{}
\titleformat{\subsubsection}{\normalsize\bfseries\color{iscte}}{\thesubsubsection}{0.8em}{}

\pagestyle{fancy}
\fancyhf{}
\fancyhead[L]{\small\color{cinza}\TituloCurto}
\fancyhead[R]{\small\color{cinza}SRSI --- Projeto, Parte 1}
\fancyfoot[C]{\small\thepage}
\renewcommand{\headrulewidth}{0.4pt}

% Nomes em português de Portugal
\addto\captionsportuguese{%
  \renewcommand{\contentsname}{Índice}%
  \renewcommand{\listfigurename}{Lista de Figuras}%
  \renewcommand{\listtablename}{Lista de Tabelas}%
  \renewcommand{\refname}{Referências}%
  \renewcommand{\abstractname}{Resumo}%
  \renewcommand{\tablename}{Tabela}%
  \renewcommand{\figurename}{Figura}%
  \renewcommand{\appendixname}{Anexo}%
}

% ---------------------------------------------------------------------------
% Macros
% ---------------------------------------------------------------------------
% Marcas usadas nas tabelas comparativas
\newcommand{\Sim}{\textcolor{bom}{\ding{51}}}          % abordado
\newcommand{\Parc}{\textcolor{parcial}{\ding{108}}}    % abordado parcialmente
\newcommand{\Nr}{\textcolor{cinza}{--}}                % não reportado
\newcommand{\Nao}{\textcolor{mau}{\ding{55}}}          % ausente por desenho

% Identificadores dos estudos incluídos (E01..E27) com citação
\newcommand{\E}[1]{\textsf{E#1}}

% Colunas de tabela
\newcolumntype{L}[1]{>{\raggedright\arraybackslash}p{#1}}
\newcolumntype{C}[1]{>{\centering\arraybackslash}p{#1}}
\newcolumntype{Y}{>{\raggedright\arraybackslash}X}

% Caixas de definição/observação
\newcommand{\caixa}[2]{%
  \begin{center}\begin{tikzpicture}
    \node[draw=acento, fill=fundotab, rounded corners=2pt, inner sep=8pt,
          text width=0.9\linewidth, align=justify] {\textbf{#1.}\ #2};
  \end{tikzpicture}\end{center}}
